\section{Consideraciones éticas}
\label{sec:etica}

El trabajo con datos generados por personas exige un cuidado adicional. Este proyecto
asume expl\'\i{}citamente un marco \'etico basado en la minimizaci\'on, la anonimizaci\'on,
la transparencia y la interpretaci\'on prudente.

\subsection{Finalidad y consentimiento}

El an\'alisis se realiza con fines exclusivamente acad\'emicos y sobre agregados. No
persigue evaluar individuos, vigilar conductas ni tomar decisiones que afecten a
personas. Los participantes del grupo no fueron informados individualmente del estudio;
esta circunstancia, com\'un en la investigaci\'on de datos observacionales, impone una
obligaci\'on reforzada de no da\~nar y de limitar el an\'alisis a lo estrictamente
necesario. La literatura sobre \'etica en investigaci\'on con redes sociales
\citep{townsend2016social,freelon2019computational} subraya precisamente este punto:
la ausencia de consentimiento expl\'\i{}cito no exime, sino que acrecienta, la
responsabilidad del investigador.

\subsection{Anonimización y minimizaci\'on}

Los identificadores de emisor se transforman en \emph{hashes} en todas las salidas. El
informe no reproduce mensajes completos atribuibles a personas y se limita a ejemplos
genéricos o construidos. El corpus bruto no se versiona en el repositorio p\'ublico y se
mantiene fuera de \'el mediante reglas de exclusi\'on. Se aplica el principio de
minimizaci\'on: solo se conservan las variables necesarias para el an\'alisis.

\subsection{Interpretaci\'on prudente}

Una etiqueta emocional autom\'atica no es un diagn\'ostico ni una afirmaci\'on sobre el
estado real de una persona. El avatar emocional del panel interactivo y todas las
visualizaciones representan agregados del modelo, no mediciones psicol\'ogicas. Esta
distinci\'on se documenta de forma visible en las herramientas y se reitera en el
informe para prevenir interpretaciones abusivas.

\subsection{Transparencia}

Se documentan p\'ublicamente el m\'etodo, los par\'ametros, las versiones del diccionario
y del software, las limitaciones y los sesgos conocidos. La reproducibilidad
metodológica no es solo una virtud t\'ecnica, sino una
condici\'on \'etica: permite que terceros auditen y cuestionen los resultados.

\subsection{Riesgos y mitigaciones}

\begin{itemize}[leftmargin=1.4em]
  \item \textbf{Reidentificaci\'on.} Mitigada mediante el uso exclusivo de hashes y la
        no publicaci\'on de mensajes atribuibles.
  \item \textbf{Uso indebido de las etiquetas.} Mitigada mediante advertencias expl\'\i{}citas
        sobre el carácter aproximado y agregado de las etiquetas.
  \item \textbf{Sesgo del diccionario.} Mitigada mediante documentaci\'on abierta y la
        l\'\i{}nea de trabajo futuro dedicada a su ampliaci\'on validada.
  \item \textbf{Vigilancia.} Mitigada mediante el an\'alisis exclusivamente agregado y la
        ausencia de perfilado individual por defecto.
\end{itemize}

\subsection{Cumplimiento normativo}

El tratamiento se alinea con los principios de las normas de protecci\'on de datos
aplicables (finalidad, minimizaci\'on, limitaci\'on del plazo de conservaci\'on y
seguridad), entendiendo que el uso es de investigaci\'on acad\'emica y con datos
anonimizados. Cualquier publicaci\'on derivada respetar\'a estas mismas garant\'\i{}as.
